
Machine learning practitioners face method selection challenges with limited systematic guidance. This work investigates whether meta-learning can predict optimal method families (Linear, Polynomial, RNN, Augmentation) from dataset characteristics by training a Random Forest classifier on published benchmark results. The approach decomposes into three testable sub-hypotheses: data collection feasibility (h-e1), feature-method correlations (h-m1), and meta-classifier training (h-m2). Experiments collected 29 benchmarks from accessible sources (OGB, FedML, GitHub repositories, published papers), but literature mining yielded sparse metadata: 13.8\% coverage for sample size, 0\% for dimensionality, and zero variance in class imbalance despite 75.9\% nominal coverage. This sparsity prevented correlation testing (h-m1: zero computable correlations due to insufficient feature diversity) and caused meta-classifier failure (h-m2: 25.6\% accuracy with one usable feature, below the 48.3\% majority baseline). The negative results stem from data limitations rather than hypothesis refutation. Literature mining successfully identifies benchmark sources but cannot extract dataset characteristics---APIs provide identifiers without statistics, READMEs document usage without sample sizes, and papers report method rankings without dataset properties. A two-stage collection protocol is required: source identification (achieved) and characteristic extraction via dataset downloads (the bottleneck). This work quantifies the metadata gap, proposes community-driven benchmark metadata repositories, and reframes the research question from testing whether meta-learning works to identifying the infrastructure needed for fair evaluation.
